\section{Dataset Analysis}
\label{sec:analysis}

\subsection{Hypothesis}
\label{sec:hypothesis}

The new Elixir typechecker is designed to be more error-prone than Dialyzer.
% Dialyzer is built around success typing: it warns only when it can prove that a discrepancy is certain, accepting many function definitions whose behaviours are merely suspicious.
The set-theoretic typechecker performs a structural analysis of each function body against its declared and inferred types and warns on definitions that violate them.
On this setting, every function definition Dialyzer warns should also be warned by the typechecker, and the typechecker should additionally find potential type errors Dialyzer cannot detect.
Formally:
\begin{gather*}
  \text{Dialyzer fails} \;\Rightarrow\; \text{Typechecker fails} \\
  \text{Dialyzer passes} \;\Rightarrow\; \text{Typechecker fails $\vee$ passes}
\end{gather*}
We expect the implications to hold together with the strict superset expectation that the typechecker fails on substantially more definitions than Dialyzer.

To judge whether the hypothesis holds we examine the \emph{intersection} of the two tools' warnings, that is the fraction of both warned by Dialyzer and the typechecker over only Dialyzer-warned:
\begin{gather*}
  \frac{\text{Dialyzer fails} \;\wedge\; \text{Typechecker fails}}{\text{Dialyzer fails}} \approx 1.0
\end{gather*}
The closer the fraction to one confirms the hypothesis: the new system preserves Dialyzer's discrimination and exceeds it.
Where the fraction yields less than one, we then reason about the \emph{residual}, i.e., the entries Dialyzer warns on but the typechecker does not.

The analysis addresses three questions:
\begin{enumerate}
    \item To what extent is the typechecker stricter than Dialyzer, and how much of Dialyzer's discrimination does it subsume?
    \item Where the subsumption breaks down, why does Dialyzer catch an error that the typechecker does not, and is that a fundamental disagreement or a limitation of the present state of either the typechecker or the translation?
    \item What does the new typechecker detect that Dialyzer does not, and what does this imply for downstream use?
\end{enumerate}

\subsection{Overall Results and Verdict}
\label{sec:overall-results}

The dataset construction described in Section~\ref{sec:dataset-construction} produced \num{23073} translated function specifications.
Dialyzer analysed \num{22982} of them and failed on 91; the new typechecker analysed \num{22621} and could not analyse the remaining 452, whose host projects fail to compile under the custom compiler (a dependency that either tool rejects aborts the whole project; see Section~\ref{sec:dataset-construction}).
The 91 Dialyzer could not analyse are a subset of those 452 --- entries neither tool could analyse --- so both tools return a verdict on the same \num{22621} entries.
The cross-comparison below is therefore over those \num{22621} entries both tools resolved, with the 452 reported separately.
Table~\ref{tab:overall-results} summarises the outcomes:

\begin{table}[h]
    \centering
    \begin{tabular}{lrr}
        \toprule
        \textbf{Outcome}             & \textbf{Count} & \textbf{\% of total} \\
        \midrule
        Dialyzer passed              & \num{22685} & 98.3\% \\
        Dialyzer failed              & \num{297}   & 1.3\% \\
        Dialyzer null                & \num{91}    & 0.4\% \\
        \midrule
        Typechecker passed           & \num{18786} & 81.4\% \\
        Typechecker failed           & \num{3835}  & 16.6\% \\
        Typechecker null             & \num{452}   & 2.0\% \\
        \midrule
        Both tools passed            & \num{18688} & 81.0\% \\
        Both tools failed            & \num{199}   & 0.9\% \\
        Both tools null              & \num{91}    & 0.4\% \\
        \midrule
        Typechecker-only failed      & \num{3636}  & 15.8\% \\
        Dialyzer-only failed         & \num{98}    & 0.4\% \\
        \bottomrule
    \end{tabular}
    \caption{Tool outcomes across \num{23073} translated function specifications. A \emph{null} outcome means the tool could not analyse the entry because its host project failed to compile --- for the typechecker, under the custom development compiler. Dialyzer's 91 nulls are a subset of the typechecker's 452, the entries neither tool could analyse; percentages are of the \num{23073} total.}
    \label{tab:overall-results}
\end{table}

Two observations stand out.
First, the typechecker is decisively the stricter of the two analyses: it rejects \num{3835} entries (16.6\% of the corpus) where Dialyzer rejects only 297 (1.3\%), a thirteen-fold difference.
The set-theoretic typechecker's pass rate (81.4\%) sits below Dialyzer's (98.3\%).
Second, where the both tools fail, they overlap substantially but not completely.

\paragraph{Subsumption rate on failures.}
Dialyzer warned on 297 entries in total; the typechecker also warned on 199 of them, giving a subsumption rate on failures of $199/297 \approx 67.0\%$.
The remaining 98 Dialyzer-warned entries (the Dialyzer-only set) are the cases where the first implication appears to break down.
However, as Subsection~\ref{sec:dialyzer-only} establishes, none of these 98 represents a case in which the typechecker disagrees with Dialyzer; each is attributable to an analysis outside its scope (whole-program reachability or cross-module resolution), to a feature (\texttt{@opaque}-type enforcement) it does not yet model, or to precision lost during approximation.

\paragraph{Verdict.}
The results support the hypothesis' subsumption clause but do not qualify its strict-superset reading.
The new typechecker does warn on far more definitions than Dialyzer, just as predicted; but, as Subsection~\ref{sec:typecheck-only} shows, most of that excess is translation residuals --- struct patterns the translator could not fully resolve from the AST and positions widened to \texttt{dynamic()} --- rather than superior error detection.
The first implication holds: the typechecker reproduces 67.0\% of Dialyzer's warnings directly, and none of the residual 98 is a genuine type-theoretic disagreement (Subsection~\ref{sec:dialyzer-only}).
In short, the two analyses are complementary rather than nested: where both resolve the same types the typechecker reproduces Dialyzer's findings, but each also flags a class the other cannot, namely Dialyzer the whole-program properties and intra-project calls the typechecker leaves as \texttt{dynamic()}, and the typechecker a core of purely-static body-level return leaks that Dialyzer's optimism accepts (Subsection~\ref{sec:typecheck-only}).

\subsection{Approximation via \texttt{dynamic()}}
\label{sec:approximation-via-dynamic}

Out of the \num{23073} entries, \textbf{\num{13910} (60.3\%)} contain at least one \texttt{dynamic()} sub-expression in the translated annotation.
This widening is the single most consequential approximation in the translation, and it is the mechanism behind a large share of the residual disagreement analysed below.

When the definition of a remote type \texttt{Module.type\_name()} is unavailable at translation time---the module lives in a dependency outside the translation path---the translator conservatively replaces the reference with \texttt{dynamic()}.
Because \texttt{dynamic()} is compatible with every type (the ``dynamic mode'' of ``\nameref{sec:gradual-type-system}''), the annotation is no stronger than the original specification at its unresolved positions: checking a body against a \texttt{dynamic()} argument, the typechecker cannot derive the contradiction Dialyzer derives from the original concrete contract.
The system always hoists \texttt{dynamic()} to the root of its enclosing container~\cite{elixir-set-theoretic-types}, rewriting a nested \texttt{\{:ok, dynamic()\}} to the \emph{equal} type \texttt{dynamic(\{:ok, term()\})}.
This normalises the surface form without coarsening the check: only the widened leaf is masked, while the tag \texttt{:ok} is still required, so \texttt{dynamic(\{:error, term()\})} does not match \texttt{\{:ok, integer()\}}.
The masking is thus \emph{leaf-granular}---a \texttt{dynamic()} suppresses a mismatch only at the position it occupies, leaving every concrete position around it fully checked.

How coarse the approximation is therefore depends on \emph{what} the translator could not resolve. Two cases occur in the corpus, shown here on a map field:
\begin{lstlisting}[language=Elixir, numbers=none]
# a field's value type is unknown but its key type is an atom literal: only the value is widened,
# the map stays closed, and the key set and sibling fields are still checked
%{:id => integer(), :url => dynamic()}
# hoisted internal form: the value's dynamic() lifts to the root, :url becomes term()
dynamic(%{:id => integer(), :url => term()})

# a whole field type is unknown -- its key type is unresolved too: the field type is emitted as dynamic() => dynamic(),
# whose dynamic() key expands to every key domain, opening the map and masking the field type entirely
%{dynamic() => if_set(dynamic())}
# hoisted internal form: the dynamic() key becomes term() (every key domain), the map is open
dynamic(%{term() => if_set(dynamic())})
\end{lstlisting}
In the first form, the key \texttt{:url} is a literal that needs no resolution, so closedness survives and only one value is masked; this is the common struct case, where the field names are known from \texttt{defstruct} while a field's type is out of scope.
In the second, the key itself is an unresolved type, so the emitted \texttt{dynamic()} key is expanded through \texttt{to\_domain\_keys/1} to cover every key domain as the compiler ingests the \texttt{@assert\_type\_form} annotation; the map is then no longer closed and admits any key.
A whole-field approximation is thus strictly coarser than a value-only one---it forfeits not just a value but the map's shape.

The converse bounds how a \emph{pass} should be read: an annotation with \emph{no} \texttt{dynamic()} is checked as a fully static type---``if the user provides their own types, and those types are not \texttt{dynamic()}, then Elixir's type system behaves as a statically typed one''~\cite{elixir-set-theoretic-types}.
Where the annotation does carry \texttt{dynamic()}, those positions are accepted trivially, so a pass certifies only the concrete positions; for a wholly \texttt{dynamic()} annotation, such as the opacity cases of Subsection~\ref{sec:dialyzer-only}, a pass is no static verdict at all.

The prevalence of \texttt{dynamic()} does not invalidate the translation---as a gradual type it is a sound approximation, deferring an unresolved position to a runtime check rather than falsifying it statically~\cite{duboc2024}---but instead quantifies how much of the corpus carries cross-module type references that cannot be refined without expanding the translation scope to the full dependency tree.
The widening cuts both ways, one direction per following subsection: it suppresses some Dialyzer-detectable errors (Subsection~\ref{sec:dialyzer-only}), yet \num{2747} of the \num{3835} typechecker-rejected entries (72\%) carry a \texttt{dynamic()} and are rejected regardless (Subsection~\ref{sec:typecheck-only})---leaf-granular masking, not a blanket suppression of the typechecker.

\subsection{Dialyzer-Only Warnings: Why Dialyzer Catches What the Typechecker Cannot \emph{Yet}}
\label{sec:dialyzer-only}

The 98 Dialyzer-only entries are the crux of the hypothesis: they are precisely the cases where Dialyzer rejects a definition the typechecker accepts.
However, establishing that none of them is a true type-theoretic counterexample is what allows the first implication to stand.
Table~\ref{tab:dialyzer-cats} breaks them down by warning category; the category names and their meanings follow Dialyzer's own classification~\cite{dialyzer-manual}, as surfaced through the \texttt{dialyxir} interface used during construction~\cite{dialyxir}.
A counting caveat first: Dialyzer routinely reports several warnings for a single function definition (these 98 entries carry 134 warnings between them), whereas the new typechecker generally surfaces only one per function, a difference that matters once both tools fire on the same definition (Subsection~\ref{sec:both-warned}).
The column \textbf{N} below therefore counts the \emph{entries} that carry each category, not raw warnings.
A representative entry for each category, namely the function and the Dialyzer message it draws, none of which the typechecker reproduces, is given in the appendix ``\nameref{app:dialyzer-examples}''.

\begin{table}[h]
\centering
\footnotesize
\begin{tabularx}{\textwidth}{@{}l r X l@{}}
\toprule
\textbf{Category} & \textbf{N} & \textbf{What Dialyzer reports} & \textbf{In typechecker?} \\
\midrule
\texttt{no\_return}            & 29 & A function never returns normally; it always raises, exits, or loops. & no (whole-program) \\
\texttt{pattern\_match}        & 18 & A pattern can never match the type of the term it is matched against. & partial (local) \\
\texttt{call\_without\_opaque} & 12 & An \texttt{@opaque}-typed argument is supplied by its raw internal structure. & no (opacity) \\
\texttt{unused\_fun}           & 12 & A function is defined but never called anywhere. & no (whole-program) \\
\texttt{call\_to\_missing}     & 11 & A call targets a function missing or not exported by its module. & no (cross-module) \\
\texttt{unknown\_function}     &  9 & A call targets a function that does not exist. & no (cross-module) \\
\texttt{pattern\_match\_cov}   &  5 & A clause is already fully covered by earlier clauses and is unreachable. & partial (local) \\
\texttt{guard\_fail}           &  4 & A guard test can never succeed for the inferred types. & partial (local) \\
\texttt{call}                  &  3 & A call's arguments cannot match the callee's expected input types, so it can never succeed. & partial (local) \\
\texttt{callback\_type\_mismatch} & 1 & A \texttt{@behaviour} callback implementation's spec conflicts with the behaviour's declared callback type. & no (cross-module) \\
\texttt{opaque\_match}         &  1 & A pattern inspects the internal structure of an \texttt{@opaque} term. & no (opacity) \\
\bottomrule
\end{tabularx}
\caption{Dialyzer warning categories among the 98 Dialyzer-only entries. \textbf{N} is the number of entries carrying each category; an entry whose warnings span several categories appears in more than one row, so the column sums to 105, not 98 (7 of the 98 entries carry warnings in two categories). This remains below the 134 raw warnings quoted above, because an entry may also carry several warnings within a single category. The final column indicates whether the typechecker performs an analysis that can produce the same finding. \emph{partial (local)} marks a category the typechecker checks inside a function body and therefore reproduces when both tools resolve the same types (Subsection~\ref{sec:both-warned}), but misses here because translation widened the relevant type to \texttt{dynamic()}; \emph{no} marks a finding outside its analysis altogether, requiring \emph{whole-program} reasoning (call-graph reachability or dead code), \emph{cross-module} resolution, or an \emph{opacity} model it does not have.}
\label{tab:dialyzer-cats}
\end{table}

Classifying the 98 entries by \emph{why} the typechecker does not reproduce the warning yields a clean three-way partition:
\begin{itemize}
  \item \textbf{59 entries} fall solely in categories the set-theoretic typechecker does not analyse at all: the whole-program reachability and dead code behind \texttt{no\_return} and \texttt{unused\_fun}, and the cross-module resolution behind \texttt{unknown\_function}, \texttt{call\_to\_missing}, and \texttt{callback\_type\_mismatch}.
  \item \textbf{27 entries} carry a finding the typechecker performs \emph{locally} and reproduces elsewhere, namely an argument mismatch (\texttt{call}), a dead pattern or clause (\texttt{pattern\_match}, \texttt{pattern\_match\_cov}), or an impossible guard (\texttt{guard\_fail}), but missed here; the causes, translation widening the scrutinised type to \texttt{dynamic()} versus a genuine residual gap, are examined below.
  \item the remaining \textbf{12 entries} carry \emph{solely} an opacity warning (\texttt{call\_without\_opaque}, \texttt{opaque\_match}) for \texttt{@opaque} type that the typechecker does not model and the translation cannot represent.
\end{itemize}
None of the three is a case where the typechecker, given the same information and feature set, type-theoretically disagrees with Dialyzer. We treat each in turn.

\paragraph{Whole-program analyses outside its scope (59 entries).}
By design the new typechecker performs \emph{module type inference}~\cite{elixir-set-theoretic-types}: it infers and checks each function from the current module, the standard library, and the project's dependencies, and it treats every call to another module \emph{within the same project} as \texttt{dynamic()}.
This inference is \emph{best-effort}. It warns only when a type is \emph{certain} to fail---when \emph{every} possible typing of an expression leads to an error---so a function that passes is one it could not disprove, not one it has proven correct.
Five of Dialyzer's categories fall outside this body-local view.
\texttt{no\_return} and \texttt{unused\_fun} come from Dialyzer's \emph{whole-program} (the entire call graph at once) success-typing fixed point~\cite{lindahl2006practical}: detecting that a function never returns, or that it is never called, requires reasoning across the program's control flow and call sites---which a function-at-a-time typechecker does not do.
\texttt{unknown\_function}, \texttt{call\_to\_missing}, and \texttt{callback\_type\_mismatch} instead require \emph{cross-module} resolution (a lookup of the specific module a call names).
Since the typechecker treats every same-project call as \texttt{dynamic()}, it never looks the callee up, so it never sees that the target is missing or unexported (\texttt{unknown\_function}, \texttt{call\_to\_missing}), nor that an implementation contradicts the callback type its \texttt{@behaviour} declares (\texttt{callback\_type\_mismatch}).
On all 59 entries the typechecker likewise found no body-level fault, so it passed.
These warnings are therefore not type disagreements but findings of analyses it does not \emph{yet} perform,\footnote{The word ``yet'' is not a promise. The type system is in an early phase whose stated goal is to assess performance and collect feedback from users, and its further milestones are explicitly conditional (``if the results are satisfactory''), so these analyses are not guaranteed to be added~\cite{elixir-set-theoretic-types}.} and nothing in its theory rules them out.
Until they are added, Dialyzer stays complementary here: the 59 entries mark a difference in analytical scope, not a definition the typechecker wrongly accepts.

A further point makes the \texttt{unknown\_function} and \texttt{call\_to\_missing} warnings the weakest of all.
Dialyzer can only check a call if it has already loaded the module being called; when that module was not loaded, it simply reports the call as missing.
Such a warning says only that some called module was left out of the analysis---not that the code is wrong.
Every case here calls a module that no project loads by default: optional Erlang applications (\texttt{:eldap}), development and test tools (\texttt{IEx}, \texttt{Mix}, \texttt{ExUnit}), optional Elixir libraries (\texttt{Jason}, \texttt{PhoenixHTMLHelpers}), or the compiler's own internal helpers (\texttt{:elixir\_quote}, \texttt{Enum.\_\_in\_\_/2}).
This is why whether such a warning appears is a condition of the run, not a property of the code: the same call warns in a run where its module was not loaded for analysis and stays silent in one where it was.\footnote{Dialyzer reasons only about the modules in its lookup table (the PLT)~\cite{dialyzer-manual}, which \texttt{mix dialyzer} builds from the Erlang/OTP and Elixir core applications, the dependencies fetched for the project, and the project's own compiled code. A called module falls outside this set---and is reported as missing---when its application is not among them: an optional OTP application, a dependency scoped to \texttt{:dev} or \texttt{:test}, an optional dependency that was not fetched, or a compiler-internal module, unless the project lists it explicitly under \texttt{:plt\_add\_apps}~\cite{dialyxir}.}
The new typechecker never looks up calls to other modules in the same project, so it has no such check at all; the differing verdict here is therefore a matter of scope---the two tools are not answering the same question---rather than a type-theoretic disagreement.

\paragraph{Locally checkable but unreproduced (27 entries).}
The remaining non-opacity categories are all within the typechecker's body-local remit: an argument mismatch at a call site (\texttt{call}), a pattern or clause that can never match (\texttt{pattern\_match}, \texttt{pattern\_match\_cov}), and a guard that can never succeed (\texttt{guard\_fail}) are each a body-against-declaration inconsistency it is built to catch.
That it does catch them is visible in the both-warned set (Subsection~\ref{sec:both-warned}), where it independently emits the same finding on 36 of 78 \texttt{call} entries, 30 of 55 \texttt{pattern\_match}, and a few of the others.
That it did not reproduce these 27 is a matter of approximation and analysis scope, not of an inability to express the check itself.
For only 6 of them could the cause be direct \texttt{dynamic()} masking: the annotation still carries a \texttt{dynamic()} that can cover the position Dialyzer faults, so where Dialyzer reasons against the original concrete type and finds the call impossible or the pattern dead, the widened annotation presents the typechecker with a value that vacuously satisfies the call and makes every pattern reachable (Subsection~\ref{sec:approximation-via-dynamic}); with the type translated concretely the typechecker would be expected to reproduce the finding.
The other 21 carry \emph{no} \texttt{dynamic()} at all, and their annotations are faithful renderings (a source \texttt{map()} becomes \texttt{open\_map()}, an \texttt{any()} becomes \texttt{term()}, with nothing widened), so they are a genuine residual gap: real dead-clause, unreachable-pattern, and impossible-guard findings the typechecker simply does not make.
Each of the 21 needs whole-program success typing or compile-time constant-folding that the module-local, best-effort typechecker does not perform.
The unreachability comes from one of three sources: the return type of an intra-project call, which module-local inference assumes is \texttt{dynamic()}---the Mobilizon and Akkoma federation fetchers (functions that retrieve remote ActivityPub objects), for instance, return an error tuple where the caller expects a map; a caller-side narrowing of an argument below its declared type (\texttt{ttl\_to\_seconds/1} declares a time unit of \texttt{:second}, \texttt{:minute}, \texttt{:hour}, or \texttt{:week} but is only ever called with \texttt{:hour} or \texttt{:week}, so its \texttt{:second} and \texttt{:minute} clauses never match); or a comparison of compile-time constants it does not fold (\texttt{@env} is the build environment, \texttt{:dev} here, so the \texttt{:test} branch of \texttt{if @env == :test} is statically dead).
The Dialyzer categories are ones the typechecker can express; these instances just call for reasoning the module-local analysis does not perform---largely the same reach that yields the 59 above, surfacing under a locally-checkable category rather than as a translation artefact.

\paragraph{Opacity warnings the typechecker does not model (12 entries).}
The remaining entries carry \texttt{opaque\_match} or \texttt{call\_without\_opaque}: Dialyzer warns that code constructs or pattern-matches an \texttt{@opaque} type through its internal structure rather than its public interface, an abstraction-boundary violation, here on library types such as \texttt{Ecto.Multi} and \texttt{Mint.HTTP.t()}, and even a project-local opaque (\texttt{Concentrate.TripDescriptor}).
This is not a type error in the ordinary sense but an opacity-enforcement finding, and the new type system does not distinguish an opaque type from its underlying structure: it does not model the \texttt{@opaque} (including \texttt{@typep}) type at all.
The translation cannot supply what the target system lacks: an opaque type is bluntly translated to its structure, or, when its defining dependency is unavailable, to \texttt{dynamic()}, hence the boundary of type security levels is erased on both sides.
Empirically the widening applies to eleven of the twelve: they carry a \texttt{dynamic()} in the translated annotation and none names the opaque type at all.
Crucially, and unlike the locally-checkable group, this is \emph{not} a precision loss that resolving the dependency would repair: even with the opaque type fully available, the typechecker would still not flag the violation, because it has no notion of opacity to violate.
Tellingly, the one entry whose annotation contains no \texttt{dynamic()} at all (\texttt{ExCubicIngestion.ProcessIngestion.ingest/1}, on \texttt{Ecto.Multi}) is missed all the same---confirming that opacity, not approximation, is the operative cause.

\paragraph{Project distribution.}
The Dialyzer-only warnings are concentrated: Mobilizon contributes 50 of the 98 and Akkoma 20, with the remaining 28 spread across twelve further projects (MBTA seven; Analytics five; Absinthe and Cachex three each; Cldr and Livebook two each; and six projects with one apiece).
The concentration tracks \texttt{dynamic()} density: the projects with the most unresolved remote types are those whose Dialyzer findings the translation most often masks with the gradual type.

\subsection{Typechecker-Only Warnings}
\label{sec:typecheck-only}

\num{3636} entries triggered a warning from the typechecker that Dialyzer did not report.
At face value this is the strict-superset clause of the hypothesis; on inspection, however, most of it is not superior error detection but translation residuals and the conservatism of module-local inference.
The findings divide by kind into unreachable clauses or patterns (\num{1677}), argument mismatches at a call site (914), return-type mismatches (691), map-key accesses (217), and a residue of other checks (137).

\paragraph{Unreachable-clause findings (\num{1677}).}
The largest group is not the typechecker over-warning but its \emph{correct} reasoning from a \emph{wrong} translated type: in \num{1308} of the \num{1677}, a struct pattern is matched against a translated struct type that is \emph{disjoint} from the real struct, so the pattern can never match and the clause is reported unreachable. Each is a residue of translating structs from the AST without a compiler-level lookup table (Subsection~\ref{sec:struct-representation}).
In 752 the translated struct type lists only some of the struct's fields, so the closed map excludes valid structs and a \texttt{\%Mod\{\}} pattern cannot match; in 556 the translator could not expand an aliased struct name to its full module path, so the asserted \texttt{\_\_struct\_\_} atom names a different module than the pattern's.
A verified instance is \texttt{Scenic.Components.button/3}, whose clause head pattern-matches the struct:
\begin{lstlisting}[language=Elixir, firstnumber=last]
def button(%Graph{} = g, data, options) do
  add_to_graph(g, Component.Button, data, options)
end
\end{lstlisting}
The compiler resolves \texttt{\%Graph\{\}} to \texttt{\%Scenic.Graph\{\}} (\texttt{:"Elixir.Scenic.Graph"}), disjoint from the translated argument's \texttt{:"Elixir.Graph"} (because the translator could not expand the alias to its full module path), so the typechecker reports the clause unreachable; Dialyzer, reading the original spec, accepts the same clause, confirming the code is correct and the fault is the translation.
The warning is thus \emph{valid} --- a sound deduction from the type it was handed --- yet it exposes a translation defect rather than a program fault, so it is neither a type-theoretic disagreement nor genuine error detection; both the field-set and alias residuals would be removed by resolving structs against their defining modules.
Whether the defect surfaces at all, however, depends on the body: only a struct the body \emph{exercises}---through a pattern or a field access---is caught. The same \texttt{:"Elixir.Graph"} defect in \texttt{Scenic.Primitives.rrect/3}, which merely forwards its argument, produces no warning and is admitted into the both-pass pool, sound but unfaithful to the specification:
\begin{lstlisting}[language=Elixir, firstnumber=last]
def rrect(graph_or_primitive, rrect, opts \\ []) do
  rounded_rectangle(graph_or_primitive, rrect, opts)
end
\end{lstlisting}
The \num{1308} counted here are therefore only the \emph{exercised} residuals; an unknown further number pass silently into the pool.
The remaining 369 are non-struct: genuine dead code the typechecker surfaces and Dialyzer accepts, each a branch the program can never reach. A representative case is \texttt{predictions\_from\_entites/3}, whose \texttt{@spec} types its \texttt{env} argument as \texttt{:prod} or \texttt{:dev\_green} while the body's \texttt{case} still has a third clause:
\begin{lstlisting}[language=Elixir, firstnumber=last]
@spec predictions_from_entites([map()], integer(), :prod | :dev_green) :: [map()]
def predictions_from_entites(entities, timestamp, env) do
  # ...
  case env do
    :prod -> "prod"
    :dev_green -> "dev-green"
    :dev_blue -> "dev-blue"
  end
  # ...
end
\end{lstlisting}
Since \texttt{env} can only be \texttt{:prod} or \texttt{:dev\_green}, the \texttt{:dev\_blue} clause can never match: the typechecker reports it unreachable, while Dialyzer accepts it. The others are of the same kind---a catch-all clause the earlier clauses already make unreachable by covering every case, the right-hand side of an \texttt{||} that never runs because the left side is always truthy, or a clause whose guard can never succeed.
For a codebase moving off Dialyzer onto the new checker, this reachability and exhaustiveness feedback is precisely the kind of static signal the transition is meant to add.

\paragraph{Argument mismatches (914).}
Most are the conservatism of module-local inference: when the faulted argument is inferred \texttt{term()} or \texttt{dynamic()}---from an intra-project call, or from a broad declared type such as \texttt{map()} or \texttt{list()}---the typechecker cannot establish the callee's expected input and warns on correct code.
\texttt{Scenic.PubSub.do\_subscribe/3} is representative: its \texttt{state} argument is declared \texttt{map()}, an open map that carries no per-field types, so the \texttt{subs\_id} destructured from it is inferred \texttt{term()}:
\begin{lstlisting}[language=Elixir, firstnumber=last]
@spec do_subscribe(pid :: GenServer.server(), source_id :: atom, state :: map) :: any
def do_subscribe(pid, source_id, %{subs_id: subs_id, subs_pid: subs_pid} = state) do
  # ...
  Map.get(subs_id, source_id, [])
  # ...
end
\end{lstlisting}
\texttt{Map.get/3} expects a \texttt{map()} as its first argument, but \texttt{subs\_id} is \texttt{term()}, so the call is warned---even though \texttt{state} carries a map there at runtime, which is why Dialyzer accepts it.
The rest fall into two smaller groups. Some pass a struct whose translated type is missing a field the callee accesses, so the argument does not satisfy what the callee expects---the unrecovered-field-set residual of Subsection~\ref{sec:struct-representation}. Others call through an Elixir protocol, such as \texttt{Ecto.Queryable}, whose dispatch the checker does not model, so it cannot tell what type the protocol's implementation expects.

\paragraph{Return-type findings (691).}
These classify by the typechecker's inferred body type against the declared type into three groups.
In \num{390} the annotation itself carries a \texttt{dynamic()}: after root-hoisting the declared type is gradual, so neither a pass nor a warning here is a static verdict (Subsection~\ref{sec:approximation-via-dynamic}), and the group collects translation artefacts. \texttt{Absinthe.Phase.Debug.run/2} is representative:
\begin{lstlisting}[language=Elixir, firstnumber=last]
@spec run(any, Keyword.t()) :: {:ok, Absinthe.Blueprint.t()}
def run(input, _options \\ []) do
  Logger.debug("[Absinthe Blueprint]", input)
  {:ok, input}
end
\end{lstlisting}
\texttt{Blueprint.t()} is a struct the translation could not fully resolve, so the asserted return became \texttt{dynamic(\{:ok, \%Blueprint\{\dots\}\})}. The body returns \texttt{\{:ok, input\}} with \texttt{input} typed \texttt{term()} (it is declared \texttt{any}), so the typechecker flags the mismatch---but on a gradual type, an artefact of the translation rather than a static verdict.
In a further 125 the declared type is concrete but inference cannot pin the body's return to it and types it as the top type \texttt{term()}; since \texttt{term()} is not a subtype of a concrete declaration, the typechecker warns, whereas Dialyzer's whole-program success typing resolves the value and confirms the contract:
\begin{lstlisting}[language=Elixir, firstnumber=last]
@spec backend() :: backend()
def backend do
  get(:backend)
end
\end{lstlisting}
The \texttt{@spec} promises the return type \texttt{backend()}---defined as just the two atoms \texttt{:exla} or \texttt{:nx\_binary}---but the body calls the same module's \texttt{get/1}, whose own declared return is \texttt{term()} (its body is \texttt{Map.get(get(), key)}). \texttt{backend/0} is therefore typed \texttt{term()}, and since \texttt{term()} is not a subtype of those two atoms, the typechecker cannot show the body stays within the declared type and warns; Dialyzer accepts, because \texttt{term()} does not contradict \texttt{:exla} or \texttt{:nx\_binary} on any path.

The subtlety is why the body is \texttt{term()} rather than \texttt{dynamic()}, since both denote ``any value''. Module type inference assumes \texttt{dynamic()} only for calls into \emph{other} modules of the same project; a call to a function in the \emph{same} module, like \texttt{get/1}, is resolved to its actual type---here \texttt{term()}~\cite{elixir-set-theoretic-types}. This is decisive, because the two tops behave oppositely against a concrete declaration: \texttt{term()} is the \emph{static} top and is not a subtype of a concrete type, so it warns, whereas \texttt{dynamic()} is the \emph{gradual} top, compatible with every type, so it never does. Had \texttt{get/1} lived in another project module its result would have been assumed \texttt{dynamic()}, and \texttt{backend/0} would have passed.
The remaining 176 are purely static---no \texttt{dynamic()} on either side---so a disagreement is a true static verdict: the body provably returns a value outside the declared type that Dialyzer's optimism accepts:
\begin{lstlisting}[language=Elixir, firstnumber=last]
@spec bounds(vectors :: nil | list(Math.vector_2())) ::
        {left :: number, top :: number, right :: number, bottom :: number}
def bounds(nil) do
  nil
end
def bounds([]) do
  nil
end
def bounds([{x, y} | vectors]) when is_list(vectors) do
  # ... reduces to a {left, top, right, bottom} tuple
end
\end{lstlisting}
For a \texttt{nil} or empty input \texttt{bounds/1} returns \texttt{nil}, outside the declared coordinate tuple; Dialyzer's optimism accepts it because another clause can return the tuple.
These purely-static leaks are the legitimate content of the strict-superset clause; representative cases are given in the appendix ``\nameref{app:typecheck-examples}''.

\paragraph{Interpretation.}
The headline ``\num{3636} warnings Dialyzer misses'' overstates the typechecker's advantage --- not because the typechecker is wrong (its warnings are valid) but because most are correct responses to a faulty \emph{input}, the translated type, rather than detections of a faulty \emph{program}.
The unreachable-clause group is dominated by struct-representation residuals of AST-based translation (\num{1308} of \num{1677}), where the checker rejects a disjoint struct type, and the return-type group by gradual (\texttt{dynamic()}-carrying) and cannot-verify cases that carry no static verdict (515 of 691).
What the typechecker alone genuinely catches is the 176 purely-static return leaks Dialyzer's optimism accepts and the 369 non-struct unreachable-clause findings its reachability checking surfaces.
The two analyses are thus better described as \emph{complementary} than nested for now: Dialyzer flags the whole-program properties and intra-project calls the typechecker leaves as \texttt{dynamic()}, and the typechecker both the return leaks Dialyzer's optimism accepts and the dead code its reachability checking surfaces.
As the migration target shifts from Dialyzer to the new checker, these two classes, the return leaks and the reachability findings, are what the codebase stands to gain.

\subsection{Entries Flagged by Both Tools}
\label{sec:both-warned}

199 entries triggered warnings from both tools, the core of the subsumption rate.
Their Dialyzer categories are led by \texttt{call} (78), \texttt{no\_return} (70), \texttt{pattern\_match} (55), and \texttt{call\_without\_opaque} (24), a mix of the locally-checkable and whole-program categories of Subsection~\ref{sec:dialyzer-only}, but here arising in functions whose bodies \emph{also} violate their declared return or argument types.
These are the unambiguous errors: a definition wrong enough that Dialyzer's success typing and the typechecker's body check independently reject it.
The typechecker reports a single warning per function (158 of the 199 carry exactly one), and which one it is depends on the category: for the locally-checkable categories it reproduces Dialyzer's finding directly (incompatible arguments on 36 of 78 \texttt{call} entries, a dead pattern on 30 of 55 \texttt{pattern\_match}), whereas for the whole-program categories it cannot, and rejects instead on a coexisting body-level fault; across the 35 entries carrying only whole-program or cross-module Dialyzer categories, that fault is an unreachable clause in 21, a map-key access in 6, an argument mismatch in 5, a return-type mismatch in 2, and another check in 1.
Either way the entry is co-rejected: the typechecker need not match Dialyzer's diagnosis, only fail the function on some ground.
Their existence confirms that the two analyses agree wherever a definition's fault is visible to both, and that the disagreement studied above is confined to the analytical-scope and approximation boundaries, not to the verdict on plainly faulty code.
One entry per Dialyzer category from this set, paired with both tools' messages, is given in the appendix ``\nameref{app:both-examples}''.

\subsection{Implications for Downstream Use}
\label{sec:downstream-implications}

The cross-validation conducted by Dialyzer and Elixir's new typechecker is not only a quality check on the translation; it directly shapes how the resulting dataset is consumed in the LLM training pipeline in the next part, whose goal is to infer a function type that the typechecker accepts.
Three observations transfer directly.

\paragraph{Keeping only typechecker-verified labels.}
The models are trained to produce annotations the typechecker accepts, so every training label must itself be an annotation the typechecker admits.
We therefore keep only the \emph{both-pass} entries, those accepted by \emph{both} Dialyzer and the new typechecker.
This is a substantive filter: it removes \num{4385} entries (19.0\% of the corpus), the great majority typechecker rejections, plus the 452 the typechecker could not analyse and the 297 Dialyzer warnings.
About a fifth of the translated annotations are thus rejected by, or unverifiable under, at least one tool (genuinely inconsistent for some, but for most only unverifiable by the typechecker's conservative inference; see Subsection~\ref{sec:typecheck-only}), and excluding them leaves only labels the typechecker positively accepts.

\paragraph{Gradual types in the labels.}
Those labels can carry \texttt{dynamic()}, inherited from the translation and mostly an artefact of unresolved remote types rather than a deliberate gradual-typing choice (Subsection~\ref{sec:approximation-via-dynamic}), so the model is trained to reproduce annotations that may include \texttt{dynamic()}.
A length filter and a subproject-stratified split (Subsection~\ref{sec:llm-data-prep}) give the final training and test sets.

\paragraph{Scoring predictions.}
Dialyzer plays no part in scoring predictions: its role is confined to building the dataset (the agreement study of this chapter and the both-pass filter).
Predictions are scored against their reference annotation with the new type system's own compatibility relation, \texttt{compatible?/2} (Section~\ref{sec:compatibility-eval}), so the typechecker-accepted references are the standard each prediction is measured against.

\subsection{Summary}

\begin{itemize}
    \item \textbf{Coverage}: all \num{23073} specs from 383 repositories were translated; Dialyzer analysed all but 91 and the new typechecker all but 452, whose host projects fail to compile under the development-branch compiler; the 91 are among the 452, leaving \num{22621} with a verdict from both tools.
    \item \textbf{Relative strictness}: the typechecker rejects \num{3835} entries (16.6\%) against Dialyzer's 297 (1.3\%); pass rates are 81.4\% versus 98.3\%. The typechecker is decisively the stricter tool, as hypothesised.
    \item \textbf{Subsumption}: the typechecker independently warns on 67.0\% of Dialyzer's failing entries (199 of 297), reproducing Dialyzer's finding directly for the locally-checkable categories and co-rejecting the rest on a separate body-level fault. The 98 it does not co-warn are \emph{not} type-theoretic counterexamples, but split between whole-program analyses it does not perform, locally-checkable findings masked by \texttt{dynamic()}, and \texttt{@opaque} type violations it does not model (Subsection~\ref{sec:dialyzer-only}).
    \item \textbf{Approximation cost}: 60.3\% of entries contain \texttt{dynamic()}; it suppresses some Dialyzer-visible errors but leaves 72\% of typechecker rejections intact, so it is not a blanket suppression.
    \item \textbf{Typechecker-only (15.8\%)}: its dominant kind is unreachable-clause findings, of which \num{1308} of \num{1677} reject struct patterns the translator could not fully resolve from the AST (incomplete field sets or unresolved module paths) rather than genuine warnings; the return-type group is likewise dominated by gradual (\texttt{dynamic()}-carrying) and cannot-verify cases carrying no static verdict. What it alone genuinely catches is 176 purely-static return leaks Dialyzer's optimism accepts and 369 non-struct dead-clause findings its reachability checking surfaces (Subsection~\ref{sec:typecheck-only}).
    \item \textbf{Both tools (0.9\%)}: definitions faulty enough that both analyses independently reject them, confirming agreement wherever a fault is visible to both.
    \item \textbf{Training pool}: the both-pass filter yields \num{18688} typechecker-verified entries; this single pool, length-filtered and split (Subsection~\ref{sec:llm-data-prep}), is the corpus on which the next part fine-tunes and compares two model architectures, a decoder-only and an encoder--decoder model.
\end{itemize}
The hypothesis is supported in its subsumption clause but qualified in its strict-superset clause.
The set-theoretic typechecker subsumes Dialyzer wherever both reason over the same resolvable types. Its residual blind spots on Dialyzer's failures are attributable to whole-program analyses it has not \emph{yet} implemented, to features such as \texttt{@opaque}-type enforcement it does not yet model, and to translation-time approximation, not to any case where it wrongly accepts a definition Dialyzer soundly rejects.
It also warns far more often than Dialyzer.
Although that excess is mostly translation residuals --- struct patterns the translation could not fully resolve from the AST, and positions widened to \texttt{dynamic()} --- a genuine core remains that the typechecker alone catches: 176 purely-static return leaks and 369 non-struct dead-clause findings.
The two analyses are therefore best read as complementary: each detects a class the other cannot.
